\section{Discussion}\label{sec:discussion}

\subsection{Physical interpretation of results}
Our results highlight an intrinsically steady-state-centered organization of cancer--immune dynamics: in the strong-Allee scenarios considered here, relevant homogeneous states act as unstable references rather than attractors. This is reflected by Re\,$\lambda_{\max}>0$ in the local analysis; in space, the same instability manifests as pattern selection and coarsening rather than convergence to a homogeneous equilibrium. These steady states organize spatial routes, while nonreciprocity tilts the phase flow and biases the geometry of cancer--healthy--immune tracking.

Two complementary observables make this organization explicit. First, field snapshots in temporal panels (Fig.~\ref{fig:fields_panel_strong_mu1_u}) show a rapid transition from disordered initial conditions to mesoscopic domains whose morphology depends on the control protocol. Second, correlation length grids $\xi(t)$ (Fig.~\ref{fig:corr_grid_auto_cross}) quantify the \emph{emerging length scale} of these domains, separating single-field coherence from inter-field co-localization. Across scenarios, $\xi(t)$ grows subdiffusively, consistent with diffusion-limited coarsening, and is well described by $\xi(t)\sim e^{\alpha}t^{1/2}$ over the simulated horizon. The thermodynamic diagnostic in Sec.~\ref{sec:thermodynamics} complements this picture through the diffusive entropy production: $\Sigma_{\mathrm{diss}}(t)$ is largest when interfaces are sharp and decreases as the fields smooth and coarsen. Its species-resolved components identify which population field controls the remaining spatial organization.

\subsection{Role of \texorpdfstring{$\mathcal{R}$}{R} (variational and reciprocity structure)}
The parameter $\mathcal{R}$ acts as a structural switch that reorganizes the spatial spectrum: for $\mathcal{R}=1$ patterns are typically less fragmented, with smoother interfaces and fewer high-frequency features, consistent with attenuation of short-wavelength modes. This qualitative observation is supported by correlation length grids, where $\mathcal{R}$ shifts $\xi(t)$ systematically in the strong Allee scenarios (Fig.~\ref{fig:corr_grid_auto_cross}). The effect is strongest in immune-centered observables: $i$--$i$, $c$--$i$, and $s$--$i$ increase their prefactors when $\mathcal{R}=1$, indicating more coherent immune domains and stronger spatial tracking of tumor and tissue structure. Importantly, $\mathcal{R}$ does not stabilize dynamics in the tested parameter set; instead it functions as a \emph{morphological selector} that changes the geometry of transient trajectories and the characteristic scale of domains.

\subsection{Adaptive immunological control \texorpdfstring{$u(\mathbf{x},t)$}{u(x,t)} as geometric coupling}
Adaptive control introduces a spatially directed action that is stronger in regions with high $c$ and low $i$. Field panels (Fig.~\ref{fig:fields_panel_strong_mu1_u}) show that activating/deactivating $u$ changes the route by which domains form and fuse. Quantitatively, $c$--$c$ tracks coherence of tumor domains, while $c$--$i$ tracks cancer--immune co-localization, i.e., the degree to which immunity follows tumor structure. The two-column organization in Fig.~\ref{fig:corr_grid_auto_cross} separates these effects visually: autocorrelations measure domain compactness within each population, whereas cross-correlations measure the geometry of interaction between populations. In the reciprocal completion of the dynamics, equivalent to free-energy minimization, the same control has no comparable impact on diffusive entropy production; its thermodynamic signature is therefore a genuinely nonreciprocal effect. This provides a concrete and measurable bridge between mechanistic control rules and emerging spatial organization.

\subsection{Comparison with previous models}
Relative to reciprocal Lotka--Volterra models and classical immune--tumor ODE/PDE models without Allee terms, the key qualitative difference is that a \emph{cooperative Allee effect} introduces a hard/smooth threshold separating collapse from proliferation, while nonreciprocity breaks the symmetry of interaction terms and biases transient routes. In spatial configurations, this combination shifts the pattern formation landscape from purely diffusion-driven instabilities to \emph{threshold-mediated} domain formation, where Allee type establishes the ease of domain expansion (weak) or imposes a critical barrier limiting domain sizes (strong). Control protocols then act mainly by selecting spatial scale and coupling structure, rather than creating stable homogeneous coexistence.

\subsection{Implications and limitations}
The present scenario set supports a robust qualitative conclusion: within explored parameters, the system operates near unstable steady states and the biologically relevant outcome is controlled \emph{containment vs escape} rather than asymptotic stability. Stability may still emerge in extended sweeps (e.g., reduced $r_c$ or improved immune recruitment via larger $r_d,\delta$), but spatial results already show that control and Allee strength strongly constrain morphology and characteristic length scales.

Limitations follow naturally. (i) The spatial horizon is short ($T=1$ in model units), so reported scaling and morphologies should be interpreted as transient-but-organized regimes. (ii) We report one realization per scenario (nb=1), so ensemble variability is not yet quantified; the correlation length framework provides a direct route to such ensemble statistics in future work. (iii) Adaptive control may introduce numerical stiffness in some runs (diagnostic triggers), suggesting that solver tolerances and control smoothness must be co-designed to ensure robust long-time integration.
